\documentclass[12pt]{article}
\usepackage[utf8]{inputenc}
\usepackage[T1]{fontenc}
\usepackage{hyperref}
\title{Tagdal Language Resources for K-12 Curriculum Development}
\author{Compiled by Claude Sonnet 5 (claude-sonnet-5) for Andr\'as Kornai}
\date{}
\begin{document}
\maketitle

\section{Introduction}
Tagdal (ISO 639-3 tda; Glottocode tagd1238) is spoken in Niger by the Igdalen. It is usually classed as Northern Songhay, with heavy Tuareg-Berber influence, and is often described as a mixed language. Wikipedia cites a 2021 figure of about 65,000 speakers, which differs from the roughly 27,000 in the survey's reference description.

\section{Primary Educational Resources}
No primers, textbooks or school programmes for Tagdal turned up in search.

\section{Linguistic and Reference Materials}
\begin{itemize}
\item Benitez-Torres, \emph{A Grammar of Tagdal: a Northern Songhay Language} (2021, LOT series, Leiden University). Open access, and the most comprehensive source found. \url{https://scholarlypublications.universiteitleiden.nl/access/item%3A3240587/view}
\item Related Leiden repository items from the same dissertation record. \url{https://scholarlypublications.universiteitleiden.nl/access/item%3A3240581/view}
\item Wikipedia, ``Tagdal language'': an overview with references. \url{https://en.wikipedia.org/wiki/Tagdal_language}
\item LINGUIST List 33.1683: an announcement relating to the grammar. \url{https://linguistlist.org/issues/33/1683/}
\end{itemize}
No dictionary was found in the search results.

\section{Religious and Community Resources}
None found.

\section{Digital and Audio Resources}
None found.

\section{Curriculum Development Context}
The only substantial resource found is the academic grammar. A curriculum would have to start from that, from community consultation, and from new teaching material. No practical orthography or literacy material was found.
\end{document}